\documentclass[11pt,a4paper]{scrreprt}
\usepackage[T1]{fontenc}
\usepackage[utf8]{inputenc}
\usepackage[ngerman]{babel}
\usepackage{lmodern,microtype,geometry,xcolor,booktabs,tabularx,longtable,enumitem,hyperref,listings,tikz,tcolorbox,fancyhdr}
\usetikzlibrary{arrows.meta,positioning,shapes.geometric,fit,backgrounds}
\geometry{margin=24mm,headheight=14pt}
\definecolor{navy}{HTML}{17365D}\definecolor{blue}{HTML}{2F75B5}\definecolor{cyan}{HTML}{DDEBF7}\definecolor{green}{HTML}{E2F0D9}\definecolor{orange}{HTML}{FCE4D6}\definecolor{graybox}{HTML}{F2F4F7}
\hypersetup{colorlinks=true,linkcolor=navy,urlcolor=blue,pdftitle={Von Prompts zu KI-Workflows}}
\pagestyle{fancy}\fancyhf{}\lhead{Von Prompts zu KI-Workflows}\rhead{Praxisleitfaden}\cfoot{\thepage}
\setlist{nosep,leftmargin=*}
\lstset{basicstyle=\ttfamily\small,breaklines=true,backgroundcolor=\color{graybox},frame=single,rulecolor=\color{gray!30},xleftmargin=4pt,xrightmargin=4pt,literate={ä}{{\"a}}1 {ö}{{\"o}}1 {ü}{{\"u}}1 {Ä}{{\"A}}1 {Ö}{{\"O}}1 {Ü}{{\"U}}1 {ß}{{\ss}}1}
\newtcolorbox{merke}{colback=cyan,colframe=blue,title=Merksatz}
\newtcolorbox{praxis}{colback=green,colframe=green!50!black,title=Praxis}
\newtcolorbox{achtung}{colback=orange,colframe=orange!70!black,title=Achtung}
\tikzset{box/.style={draw=navy,rounded corners,fill=cyan,minimum width=30mm,minimum height=10mm,align=center},tool/.style={draw=blue,rounded corners,fill=green,minimum width=24mm,minimum height=9mm,align=center},arr/.style={-{Latex[length=2mm]},thick,draw=navy}}
\title{\Huge\bfseries Von Prompts zu KI-Workflows\\[6mm]\Large Commands, Werkzeuge, Apps, Skills, Agenten und MCP verstehen und praktisch erproben}
\author{Didaktischer Praxisleitfaden}\date{Neuausgabe 2026}
\begin{document}
\maketitle
\begin{abstract}
Dieses Buch erklärt den Wandel vom reinen Textdialog zum werkzeuggestützten KI-Arbeitsprozess. Es trennt echte Produktfunktionen von selbst definierten Kurzbefehlen, zeigt ein persönliches Command-System und führt in zehn direkt ausführbaren Experimenten von einfachen Prompts bis zu sicheren, mehrstufigen Workflows. Alle Diagramme sind mit TikZ erzeugt.
\end{abstract}
\tableofcontents

\chapter{Vom Prompt zum Arbeitsauftrag}
Ein klassischer Prompt verlangt eine Antwort. Ein moderner Arbeitsauftrag beschreibt dagegen ein Ergebnis, Randbedingungen und Qualitätskriterien. Das System kann daraus Teilschritte ableiten, Werkzeuge wählen, Zwischenergebnisse prüfen und ein nutzbares Artefakt erzeugen.
\begin{center}\begin{tikzpicture}[node distance=15mm]
\node[box] (u) {Nutzerziel}; \node[box,right=of u] (o) {Orchestrierung}; \node[box,right=of o] (e) {Ergebnis};
\draw[arr] (u)--node[above]{Auftrag}(o); \draw[arr] (o)--node[above]{geprüft}(e);
\node[tool,below=of o,xshift=-28mm] (m) {Modell};\node[tool,below=of o] (w) {Werkzeuge};\node[tool,below=of o,xshift=28mm] (d) {Daten};
\draw[arr] (o)--(m);\draw[arr] (o)--(w);\draw[arr] (o)--(d);
\end{tikzpicture}\end{center}

\begin{merke}Ein guter Auftrag nennt Ziel, Kontext, Grenzen und Erfolgskriterien. Er muss nicht jeden Mausklick vorgeben.\end{merke}
\section{Vier Ebenen eines guten Auftrags}
\begin{enumerate}\item \textbf{Ziel:} Was soll am Ende vorliegen?\item \textbf{Kontext:} Für wen und wofür?\item \textbf{Grenzen:} Budget, Zeit, Quellen, Format und Berechtigungen.\item \textbf{Prüfung:} Woran erkennt man ein gutes Ergebnis?\end{enumerate}
\begin{praxis}Vergleiche: ``Erkläre RAG'' mit ``Erstelle für einen Softwareentwickler einen 20-minütigen Lernpfad zu RAG, mit Mini-Beispiel, typischen Fehlern und Selbsttest.''\end{praxis}

\chapter{Was Commands wirklich sind}
Das Wort \emph{Command} wird für zwei Dinge benutzt. Erstens gibt es produktseitige Schnellschalter, etwa den dokumentierten Start von Deep Research über \texttt{/Deepresearch}. Zweitens gibt es frei definierte Kurzbefehle wie \texttt{/analyse}. Bei diesen ist der Schrägstrich keine Magie: Er funktioniert, weil seine Bedeutung zuvor als Anweisung vereinbart wurde.
\begin{center}\begin{tikzpicture}[node distance=12mm]
\node[box] (c) {``Command''};\node[tool,below left=of c] (p) {Produktfunktion\\Werkzeug wird gestartet};\node[tool,below right=of c] (k) {Konvention\\Prompt wird expandiert};\draw[arr](c)--(p);\draw[arr](c)--(k);
\end{tikzpicture}\end{center}

\begin{achtung}Verfügbare Menüpunkte, Namen und Limits können sich nach Tarif, Plattform, Region und Rollout unterscheiden. Die aktuelle Oberfläche ist maßgeblich.\end{achtung}

\chapter{Das persönliche Command-System}
Die folgenden Befehle bilden eine kleine, wiederverwendbare Arbeitssprache.
\begin{longtable}{@{}p{.22\textwidth}p{.7\textwidth}@{}}\toprule
\textbf{Command}&\textbf{Standardverhalten}\\\midrule\endhead
\texttt{/analyse}&Fakten, Annahmen, Ursachen, Zusammenhänge und Lösungen trennen.\\
\texttt{/entscheidung}&Kriterien, Optionen, Risiken und Folgewirkungen bewerten; Empfehlung geben.\\
\texttt{/vergleich}&Alternativen nach sinnvollen Kriterien vergleichen.\\
\texttt{/plan}&Ziel in realistische Schritte, Abhängigkeiten und nächste Aktionen zerlegen.\\
\texttt{/recherche}&Aktuelle Informationen mit belastbaren Quellen untersuchen.\\
\texttt{/kritik}&Denkfehler, Lücken, Risiken und bessere Alternativen suchen.\\
\texttt{/code}&Ausführbaren, wartbaren Code entwerfen, testen und erklären.\\
\texttt{/sysml}&Anforderungen, Struktur, Verhalten, Schnittstellen und Verifikation modellieren.\\
\texttt{/unterricht}&Lernziele, Einstieg, Erklärung, Übung und Musterlösung erstellen.\\
\texttt{/dokument}&Material logisch ordnen und als professionelles Dokument ausarbeiten.\\
\texttt{/todo}&Aufgaben, Anlass, Priorität, Frist und nächsten Schritt extrahieren.\\
\texttt{/erklaere}&Vom Einfachen zum Komplexen erklären und Begriffe einführen.\\
\texttt{/lernen}&Praxisorientierten Lernpfad mit Übungen erstellen.\\\bottomrule
\end{longtable}
\section{Modifier}
\begin{tabularx}{\textwidth}{@{}lXlX@{}}\toprule
\texttt{+kurz}&nur Wesentliches&\texttt{+tief}&mehr Zusammenhänge\\
\texttt{+quellen}&Aussagen belegen&\texttt{+schritt}&konkrete Abfolge\\
\texttt{+praxis}&Umsetzung priorisieren&\texttt{+kritisch}&Gegenargumente suchen\\
\texttt{+tabelle}&tabellarisch ausgeben&\texttt{+datei}&passendes Artefakt erzeugen\\\bottomrule
\end{tabularx}
\begin{lstlisting}
/entscheidung +kritisch +tabelle
Soll unser Team die Wissensdatenbank selbst bauen oder eine Plattform nutzen?
\end{lstlisting}

\chapter{Tool Calling und Orchestrierung}

Ein Sprachmodell erzeugt zunächst Sprache. Viele reale Aufgaben erfordern jedoch aktuelle Informationen, genaue Berechnungen, Dateizugriff oder Aktionen in anderen Systemen. Tool Calling verbindet das Modell mit solchen Fähigkeiten. Das Modell entscheidet, welches Werkzeug benötigt wird, erzeugt einen strukturierten Aufruf und verarbeitet die zurückgegebenen Daten.

\begin{center}\begin{tikzpicture}[node distance=8mm]
\node[box] (a) {Auftrag};\node[box,below=of a] (v) {Verstehen und planen};\node[box,below=of v] (t) {Werkzeug wählen};\node[tool,below left=of t] (s) {Websuche};\node[tool,below=of t] (f) {Dateien};\node[tool,below right=of t] (r) {Rechnen};\node[box,below=22mm of t] (p) {Prüfen und synthetisieren};\node[box,below=of p] (z) {Ergebnis};
\draw[arr](a)--(v);\draw[arr](v)--(t);\draw[arr](t)--(s);\draw[arr](t)--(f);\draw[arr](t)--(r);\draw[arr](s)--(p);\draw[arr](f)--(p);\draw[arr](r)--(p);\draw[arr](p)--(z);
\end{tikzpicture}\end{center}


\section{Der typische Werkzeugzyklus}

Ein Werkzeugzyklus besteht aus fünf Schritten:

\begin{enumerate}
\item \textbf{Bedarf erkennen:} Welche Information oder Handlung fehlt?
\item \textbf{Werkzeug wählen:} Suche, Rechner, Dateiwerkzeug, Browser oder App.
\item \textbf{Aufruf formulieren:} Parameter wie Suchbegriff, Dateiname oder Zeitraum bestimmen.
\item \textbf{Ergebnis interpretieren:} Rückgabe auf Relevanz, Fehler und Vollständigkeit prüfen.
\item \textbf{Weiterarbeiten oder abschließen:} Nächsten Schritt wählen oder Ergebnis ausgeben.
\end{enumerate}

Ein einzelner Auftrag kann mehrere Zyklen enthalten. Für einen Produktvergleich müssen beispielsweise zuerst aktuelle Daten recherchiert, anschließend Werte normalisiert, Kosten berechnet und schließlich Empfehlungen formuliert werden.

\section{Was Orchestrierung bedeutet}

Orchestrierung ist die Koordination dieser Schritte. Sie verwaltet Reihenfolge, Abhängigkeiten, Zwischenergebnisse und Abbruchbedingungen. Gute Orchestrierung berücksichtigt, dass manche Schritte parallel möglich sind, andere aber ein vorheriges Ergebnis benötigen.

\begin{lstlisting}
Ziel verstehen
  -> Quellen recherchieren
  -> Daten vereinheitlichen
  -> Berechnung durchführen
  -> Widersprüche prüfen
  -> Bericht erzeugen
  -> Ausgabe visuell kontrollieren
\end{lstlisting}

\section{Warum Prüfschritte wichtig sind}

Ein technisch erfolgreicher Werkzeugaufruf ist noch kein inhaltlich richtiges Ergebnis. Eine Websuche kann irrelevante Treffer liefern, eine Datei kann veraltet sein, eine Tabelle kann Einheiten vermischen und ein Browser kann eine falsche Schaltfläche auswählen. Deshalb sollte der Workflow mindestens folgende Fragen beantworten:

\begin{itemize}
\item Stammt das Ergebnis aus der erwarteten Quelle und dem richtigen Zeitraum?
\item Sind Einheiten, Währungen und Bezugsgrößen konsistent?
\item Fehlen Daten, oder wurden leere Werte stillschweigend ersetzt?
\item Trägt die Quelle wirklich die daraus abgeleitete Aussage?
\item Entspricht das Endformat dem Arbeitsauftrag?
\end{itemize}

\section{Fehlerbehandlung und Abbruch}

Robuste Workflows definieren, was bei einem Fehler geschieht. Ein fehlgeschlagener Aufruf kann einmal mit korrigierten Parametern wiederholt werden. Danach sollte entweder ein alternatives Werkzeug verwendet oder die Unsicherheit an den Nutzer eskaliert werden. Endlose Wiederholungen sind kein Zeichen von Autonomie, sondern von fehlenden Abbruchkriterien.

\begin{praxis}
Formuliere einen Auftrag, der Recherche und Berechnung kombiniert. Lass das System vor Beginn einen kurzen Werkzeugplan nennen. Prüfe nach Abschluss, ob jeder Werkzeugaufruf einen erkennbaren Beitrag zum Ergebnis geleistet hat.
\end{praxis}


\chapter{Search, Deep Research und Quellenarbeit}

Aktuelle Recherche ist mehr als das Sammeln von Links. Sie beginnt mit einer präzisen Frage, wählt passende Quellen und endet mit Aussagen, deren Herkunft und Unsicherheit nachvollziehbar sind. Je nach Aufgabe reicht eine schnelle Websuche oder es ist eine mehrstufige Tiefenrecherche erforderlich.

\section{Drei Modi der Informationsbeschaffung}

\begin{tabularx}{\textwidth}{@{}lXXX@{}}\toprule
&\textbf{Chatwissen}&\textbf{Websuche}&\textbf{Deep Research}\\\midrule
Tempo&hoch&hoch&mittel\\
Tiefe&begrenzt&punktuell&systematisch\\
Quellen&nicht zwingend&direkte Links&dokumentierter Bericht\\
Einsatz&stabile Grundlagen&aktueller Fakt&komplexe Synthese\\\bottomrule
\end{tabularx}

Chatwissen eignet sich für zeitstabile Erklärungen und erste Orientierung. Websuche ist sinnvoll, sobald sich Informationen verändert haben könnten: Preise, Gesetze, Produktmerkmale, Termine, Rollen oder Softwareversionen. Deep Research eignet sich für Fragen, bei denen viele Quellen gegeneinander abgewogen, Zeiträume eingegrenzt oder widersprüchliche Perspektiven synthetisiert werden müssen.

\section{Eine gute Forschungsfrage formulieren}

Eine offene Themenangabe wie ``KI im Unterricht'' führt leicht zu einer oberflächlichen Sammlung. Eine gute Forschungsfrage enthält Untersuchungsgegenstand, Zielgruppe, Zeitraum und Entscheidungskontext:

\begin{lstlisting}
Welche drei organisatorischen Maßnahmen reduzieren 2026 die größten
Datenschutz- und Qualitätsrisiken beim Einsatz generativer KI an deutschen
Berufsschulen, und welche belastbaren Belege gibt es für ihre Wirksamkeit?
\end{lstlisting}

\section{Quellenhierarchie und Belegqualität}

Primärquellen sind Originaldokumente: Gesetze, Herstellerdokumentation, Forschungsarbeiten, amtliche Statistiken oder direkte Bekanntmachungen. Sekundärquellen ordnen diese ein. Gute Recherche verwendet Sekundärquellen zur Orientierung, stützt zentrale Tatsachenbehauptungen aber möglichst auf Primärquellen.

Für jeden wichtigen Beleg sind vier Fragen hilfreich: Wer ist Urheber? Wann wurde die Information veröffentlicht oder aktualisiert? Welche Methode oder Datengrundlage liegt zugrunde? Belegt die Fundstelle genau die formulierte Aussage oder nur einen ähnlichen Sachverhalt?

\section{Fakten, Interpretation und Unsicherheit trennen}

Eine Quelle kann einen Messwert liefern; die Bedeutung dieses Werts ist bereits Interpretation. Deshalb sollte ein Bericht klar unterscheiden:

\begin{itemize}
\item \textbf{Beobachtung:} direkt aus einer Quelle entnommene Information,
\item \textbf{Einordnung:} Schlussfolgerung aus mehreren Beobachtungen,
\item \textbf{Annahme:} notwendige, aber nicht belegte Arbeitsgrundlage,
\item \textbf{Unsicherheit:} fehlende, widersprüchliche oder methodisch schwache Evidenz.
\end{itemize}

\begin{praxis}
Formuliere eine Forschungsfrage und lege Zeitraum, bevorzugte Primärquellen, Ausschlusskriterien und Ausgabeformat fest. Prüfe anschließend drei zentrale Aussagen direkt an den verlinkten Quellen. Notiere bei jeder Aussage, ob sie wörtlich belegt, abgeleitet oder unsicher ist.
\end{praxis}


\chapter{Apps, Skills, Plugins und MCP}

Drei Dinge werden hier oft durcheinandergebracht, dabei lassen sie sich mit einem einfachen Bild trennen: Ein \textbf{Skill} ist wie ein Kochrezept -- eine genaue Anleitung, \emph{wie} etwas gemacht wird. Eine \textbf{App} ist wie ein Hausschlüssel -- sie öffnet dir Zugang zu einem bestimmten Ort, etwa deinem Kalender oder deinen Dokumenten. Und \textbf{MCP} ist wie eine Steckdose mit europäischem Normstecker: Jedes Gerät, das den gleichen Stecker hat, passt hinein, egal von welchem Hersteller es kommt -- niemand muss für jedes neue Gerät eine neue Wand aufreißen. Moderne KI-Systeme bestehen nicht nur aus einem Modell. Sie verbinden Verhaltensregeln, Datenquellen, Werkzeuge und Benutzungsoberflächen. Die Begriffe Apps, Skills, Plugins und MCP beschreiben unterschiedliche Teile dieser Architektur.

\begin{center}\begin{tikzpicture}[node distance=12mm]
\node[box] (u) {Nutzer};\node[box,right=of u] (ai) {KI-Orchestrator};\node[tool,above right=of ai] (sk) {Skill\\Vorgehen};\node[tool,right=of ai] (ap) {App\\Zugriff};\node[tool,below right=of ai] (mcp) {MCP\\Schnittstelle};\node[box,right=22mm of ap] (x) {Externe Systeme};
\draw[arr](u)--(ai);\draw[arr](ai)--(sk);\draw[arr](ai)--(ap);\draw[arr](ai)--(mcp);\draw[arr](ap)--(x);\draw[arr](mcp)--(x);
\end{tikzpicture}\end{center}


\section{Apps: Zugriff auf externe Systeme}

Apps verbinden eine KI-Anwendung mit Diensten wie Kalendern, Dokumentablagen, Projektmanagement oder Kommunikation. Abhängig von der jeweiligen App können sie Inhalte suchen, Kontext bereitstellen, Daten vorab synchronisieren oder nach Bestätigung Aktionen ausführen. Die vorhandenen Rechte des verbundenen Kontos und die Einstellungen des Arbeitsbereichs begrenzen den Zugriff -- wie ein Hausschlüssel, der zwar deine eigene Wohnungstür öffnet, aber nicht automatisch auch die Wohnung deines Nachbarn.

Der praktische Wert liegt darin, dass Informationen nicht manuell kopiert werden müssen. Gleichzeitig entstehen neue Risiken: sensible Daten können in den Arbeitskontext gelangen, Suchergebnisse können unvollständig sein und Schreibaktionen können externe Wirkung haben.

\section{Skills: wiederverwendbare Arbeitsverfahren}

Ein Skill legt fest, \emph{wie} eine bestimmte Aufgabe bearbeitet werden soll. Er kann Anweisungen, Beispiele, Vorlagen, Prüfschritte und Hilfsprogramme enthalten. Während ein Command oft nur eine kurze Benutzerschnittstelle ist -- das Wort auf dem Rezeptumschlag --, beschreibt ein Skill den ausführlichen Ablauf dahinter -- die einzelnen Kochschritte mit Mengenangaben.

Beispiel: Ein Skill für Besprechungsvorbereitung könnte den Kalendereintrag lesen, verknüpfte Dokumente suchen, offene Entscheidungen extrahieren und daraus eine standardisierte Vorbereitung erzeugen.

\section{Einen eigenen Skill erstellen}

Genau wie du ein Familienrezept einmal sauber aufschreibst, damit es nicht jedes Mal neu aus dem Gedächtnis rekonstruiert werden muss, kannst du auch einen eigenen Skill aufschreiben. Du brauchst dafür keine Programmierkenntnisse -- ein Skill ist im Kern ein gut strukturierter Text, den du einmal sorgfältig formulierst und danach immer wieder verwendest.

\subsection{Die Bausteine eines guten Skills}

Ein brauchbarer Skill beantwortet fünf Fragen:

\begin{enumerate}
\item \textbf{Auslöser:} Wann soll der Skill verwendet werden? Ein Name, ein Command oder eine typische Formulierung genügt.
\item \textbf{Zweck:} Welches Problem löst der Skill, und für wen?
\item \textbf{Ablauf:} Welche Schritte werden in welcher Reihenfolge ausgeführt? Auch: Welche Werkzeuge -- Suche, Datei, App -- werden dabei gebraucht?
\item \textbf{Qualitätsregeln:} Woran erkennt man ein gutes Ergebnis? Was darf auf keinen Fall passieren?
\item \textbf{Beispiel:} Ein vollständiges Beispiel mit Eingabe und erwarteter Ausgabe -- wie ein Foto vom fertigen Gericht neben dem Rezept.
\end{enumerate}

\subsection{Ein Skill-Rezept zum Nachbauen}

\begin{lstlisting}
Skill: Besprechungsvorbereitung

Auslöser: "Bereite mein Gespräch mit [Name] vor"
oder der Command /briefing.

Zweck: Innerhalb weniger Minuten eine verlässliche Grundlage
für ein Gespräch liefern, ohne dass ich selbst suchen muss.

Ablauf:
1. Kalendereintrag lesen: Datum, Teilnehmer, Thema.
2. Letztes gemeinsames Protokoll oder Dokument suchen.
3. Offene Punkte und unbeantwortete Fragen daraus extrahieren.
4. Drei mögliche Gesprächsziele formulieren.

Qualitätsregeln:
- Höchstens eine halbe Seite.
- Jede Aussage aus dem Protokoll klar kennzeichnen, keine Vermutung
  als Fakt darstellen (siehe Kapitel 5).
- Wenn kein Protokoll gefunden wird, das explizit sagen statt zu raten.

Beispiel:
Eingabe: "Bereite mein Gespräch mit Frau Keller vor."
Ausgabe: Kurzbriefing mit Termin, letzten Beschlüssen,
zwei offenen Fragen und einem Gesprächsvorschlag.
\end{lstlisting}

\subsection{Wo ein Skill in der Praxis landet}

Je nach Plattform kann derselbe Skill an unterschiedlichen Stellen hinterlegt werden: in benutzerdefinierten Anweisungen für den persönlichen Gebrauch, in Projektdateien für ein Team, in einem eigens erstellten GPT oder Assistenten, oder als eigenständige, installierbare Fähigkeit, falls die Plattform ein solches Konzept anbietet. Die genaue Bezeichnung und der Speicherort ändern sich mit der Produktoberfläche -- die Struktur aus Auslöser, Zweck, Ablauf, Qualitätsregeln und Beispiel bleibt jedoch über alle Plattformen hinweg gültig.

\begin{achtung}
Ein Skill, den niemand außer dir versteht, weil er zu knapp oder zu speziell formuliert ist, wird beim nächsten Update der Oberfläche schnell nutzlos. Schreibe ihn so, dass auch ein Kollege ohne Vorwissen ihn in einer Minute verstehen würde.
\end{achtung}

\begin{praxis}
Wähle eine Aufgabe, die du mindestens einmal im Monat erledigst. Schreibe dafür einen vollständigen Skill nach dem obigen Muster: Auslöser, Zweck, Ablauf, Qualitätsregeln, Beispiel. Teste ihn zweimal mit unterschiedlichen Eingaben und ergänze eine Regel für den Fall, der beim ersten Versuch schiefging.
\end{praxis}

\section{Plugins: gebündelte Fähigkeiten}

Ein Plugin ist ein Paket, das mehrere zusammengehörige Komponenten bereitstellen kann, etwa Skills, Apps oder Vorlagen -- wie ein ganzes Kochbuch mit Rezepten, Einkaufslisten und Küchentipps in einem Band, statt jedes Blatt einzeln zu sammeln. Dadurch lässt sich eine gesamte Arbeitsdomäne gemeinsam installieren und aktualisieren. Der Begriff bezeichnet also eher die Verpackung als einen einzelnen Arbeitsschritt.

\section{MCP: gemeinsame Sprache für Werkzeuge}

Das Model Context Protocol (MCP) standardisiert, wie ein KI-Client Werkzeuge und Ressourcen eines Servers entdecken und verwenden kann -- der europäische Normstecker für KI-Werkzeuge. Ein MCP-Server kann beispielsweise eine Datenbankabfrage, eine Dokumentensuche oder eine Geschäftsfunktion anbieten. Der Client erhält Beschreibungen und strukturierte Eingabeparameter und kann passende Aufrufe erzeugen. Ohne einen solchen gemeinsamen Standard bräuchte jede KI-Anwendung für jedes einzelne Werkzeug eine eigene, individuell gebaute Verbindung -- wie früher, als jedes Land eigene Steckertypen hatte und Reisende ständig Adapter mitschleppen mussten.

MCP löst nicht automatisch Fragen der Qualität oder Berechtigung. Es schafft eine technische Schnittstelle; Authentifizierung, Rechte, Protokollierung und sichere Werkzeugdefinition bleiben Aufgaben der beteiligten Systeme -- der Normstecker verhindert nicht, dass jemand ein kaputtes Gerät einsteckt.

\begin{merke}
Skill = Arbeitsweise (das Rezept). App = verbundener Dienst (der Schlüssel). Plugin = Paket aus Fähigkeiten (das Kochbuch). MCP = standardisierte Verbindung zu Werkzeugen und Ressourcen (die Steckdose).
\end{merke}

\section{Ein vollständiges Beispiel}

Beim Auftrag ``Bereite mein nächstes Projektgespräch vor'' könnte ein Skill die benötigten Schritte festlegen. Eine Kalender-App liefert Termin und Teilnehmer, eine Dokument-App liefert das letzte Protokoll, und MCP kann die technische Schnittstelle zu einem internen Projektsystem bilden. Die Orchestrierung verbindet alles zu einem Briefing -- der Skill ist das Rezept, die Apps liefern die Zutaten aus verschiedenen Läden, und MCP sorgt dafür, dass jeder Laden dieselbe Bestellsprache versteht.

\section{Woran man den Unterschied im Alltag erkennt}

Eine einfache Eselsbrücke: Wenn du fragst ``Wie mache ich das?'', denkst du an einen Skill. Wenn du fragst ``Wo kommen die Daten her?'', denkst du an eine App. Und wenn du fragst ``Wie sprechen die Systeme technisch miteinander?'', denkst du an MCP. Ein Plugin schließlich beantwortet die Frage ``Was gehört alles zusammen und wird gemeinsam installiert?''.

\begin{praxis}
Nimm einen eigenen Workflow und ordne jedes Element einer Kategorie zu: Benutzerziel, Skill-Regel, App-Zugriff, MCP-Werkzeug, Prüfschritt und externe Aktion. Markiere, an welchen Stellen Berechtigungen notwendig sind.
\end{praxis}

\chapter{Agenten und sichere Autonomie}
Ein Agent verfolgt ein Ziel über mehrere Schritte, beobachtet Ergebnisse und passt sein Vorgehen an. Autonomie sollte abgestuft sein.
\begin{center}\begin{tikzpicture}[node distance=7mm]
\node[tool] (r) {Read\\lesen};\node[tool,right=of r] (d) {Draft\\vorbereiten};\node[tool,right=of d] (a) {Act\\ausführen};\draw[arr](r)--node[above]{mehr Wirkung}(d);\draw[arr](d)--(a);
\end{tikzpicture}\end{center}

\begin{itemize}\item \textbf{Read:} Informationen abrufen, nichts verändern.\item \textbf{Draft:} Entwurf oder Änderung vorbereiten, Nutzer entscheidet.\item \textbf{Act:} Externe Wirkung auslösen; Ziel, Umfang und Bestätigung müssen klar sein.\end{itemize}
\begin{achtung}Zahlungen, Löschungen, Veröffentlichungen, Nachrichten und Rechteänderungen sind folgenreiche Aktionen. Vorher exakte Ziele und Vorschau prüfen.\end{achtung}

\chapter{Praxislabor: zehn Experimente}
Führe jedes Experiment in einem neuen Chat aus und notiere: Auftrag, verwendete Werkzeuge, Rückfragen, Ergebnis, Fehler und Verbesserung.
\section{1. Ein Prompt, drei Qualitätsstufen}
Bitte zuerst nur um eine Erklärung. Füge dann Zielgruppe und Format hinzu. Formuliere schließlich Erfolgskriterien. Vergleiche die Ergebnisse.
\section{2. Eigener Command}
Definiere \texttt{/brief}: Ziel, Ausgangslage, drei Kernaussagen, Risiken, nächste Schritte. Teste ihn an zwei Themen und verbessere die Definition.
\section{3. Modifier kombinieren}
Teste \texttt{/vergleich +kurz}, danach \texttt{+tief +kritisch +tabelle}. Beobachte, welche Modifier miteinander konkurrieren.
\section{4. Aktualität erzwingen}
Stelle eine zeitabhängige Frage. Verlange Datum, Quellen und Kennzeichnung unsicherer Aussagen. Öffne mindestens zwei Quellen.
\section{5. Deep Research}
Wähle eine Frage mit mehreren Perspektiven. Begrenze Zeitraum und Quellentypen. Prüfe den Forschungsplan vor dem Start.
\section{6. Datei als Datenquelle}
Lade ein eigenes, unkritisches Dokument hoch. Lass Aussagen extrahieren und fordere Seiten- oder Abschnittsnachweise. Suche eine erfundene Behauptung.
\section{7. Workflow mit Zwischenschritt}
Lass aus Notizen erst eine Gliederung erstellen. Genehmige sie, bevor der Text entsteht. Vergleiche das Resultat mit einer Ein-Schritt-Anfrage.
\section{8. Read--Draft--Act}
Formuliere eine Aufgabe dreimal: nur lesen, Entwurf erstellen, Aktion vorbereiten. Markiere, an welcher Stelle eine Bestätigung erforderlich ist.
\section{9. Fehler provozieren}
Gib absichtlich widersprüchliche Randbedingungen. Prüfe, ob Rückfragen kommen oder Annahmen transparent gemacht werden.
\section{10. Eigener wiederverwendbarer Workflow}
Wähle eine monatliche Aufgabe. Definiere Eingangsdaten, Schritte, Qualitätsprüfung, Freigabepunkte und Ausgabe. Teste mit einem realen Beispiel.

\chapter{Eine persönliche Workflow-Architektur aufbauen}
Beginne klein: globale Konventionen für häufige Arbeit, projektspezifische Regeln für einen Kontext und spezialisierte Assistenten nur dort, wo eigenes Wissen oder Werkzeuge notwendig sind.
\begin{center}\begin{tikzpicture}[node distance=9mm]
\node[box] (g) {Globale Commands};\node[box,below=of g] (p) {Projektregeln und Dateien};\node[box,below=of p] (s) {Spezialisierte Skills/GPTs};\node[box,below=of s] (a) {Apps und Aktionen};\draw[arr](g)--(p);\draw[arr](p)--(s);\draw[arr](s)--(a);
\end{tikzpicture}\end{center}

\section{Vorlage für benutzerdefinierte Anweisungen}
\begin{lstlisting}
Wenn meine Nachricht mit /analyse, /entscheidung, /vergleich, /plan,
/recherche, /kritik, /code, /unterricht, /dokument, /todo, /erklaere
oder /lernen beginnt, behandle das Wort als vereinbarte Kurzform für den
entsprechenden Arbeitsmodus. Modifier: +kurz, +tief, +quellen, +schritt,
+praxis, +kritisch, +tabelle, +datei. Trenne Fakten, Annahmen und
Unsicherheit. Frage nur nach, wenn eine fehlende Angabe das Ergebnis
wesentlich verändert. Vor externen oder schwer rückgängig zu machenden
Aktionen zeige Ziel und Wirkung und hole eine Bestätigung ein.
\end{lstlisting}

\chapter{Fehlerbilder und Diagnose}

Fehler in KI-Workflows entstehen selten nur an einer Stelle. Ursache kann ein unklarer Auftrag, ungeeigneter Kontext, eine schwache Quelle, ein fehlgeschlagenes Werkzeug oder eine unzureichende Prüfung sein. Statt denselben Prompt immer länger zu machen, sollte die Diagnose die fehlerhafte Ebene bestimmen.

\begin{longtable}{@{}p{.25\textwidth}p{.31\textwidth}p{.36\textwidth}@{}}\toprule
\textbf{Symptom}&\textbf{Wahrscheinliche Ursache}&\textbf{Gezielte Verbesserung}\\\midrule\endhead
Command wird ignoriert&Bedeutung fehlt im Kontext&Definition hinterlegen, Beispiel ergänzen oder Anweisung ausschreiben.\\
Antwort ist allgemein&Ziel, Zielgruppe oder Entscheidung fehlen&Gewünschtes Artefakt und Nutzungssituation nennen.\\
Antwort ist lang, aber unbrauchbar&Keine Prioritäten oder Prüfkriterien&Maximale Länge, Kernaussage und Erfolgsmaßstab festlegen.\\
Falsche Aktualität&Keine aktuelle Quelle oder kein Stichtag&Websuche, Zeitraum und Primärquellen verlangen.\\
Quelle stützt Aussage nicht&Zusammenfassung ist zu frei&Originalquelle öffnen und konkrete Fundstelle prüfen.\\
Widersprüchliche Zahlen&Einheiten oder Bezugsgrößen vermischt&Rechenweg, Einheiten und Normalisierung offenlegen.\\
Workflow dreht Schleifen&Abbruchkriterium fehlt&Maximale Versuche und Eskalationspunkt definieren.\\
Aktion wäre riskant&Rechte oder Zielbereich zu breit&Read--Draft--Act trennen und Bestätigung verwenden.\\
Datei sieht schlecht aus&Nur Inhalt, keine visuelle Prüfung&Rendern, repräsentative Seiten prüfen und iterieren.\\
Ergebnis variiert stark&Definition zu abstrakt&Positives Beispiel, Gegenbeispiel und feste Ausgabestruktur ergänzen.\\\bottomrule
\end{longtable}

\section{Eine Diagnose in fünf Fragen}

\begin{enumerate}
\item \textbf{Auftrag:} War das gewünschte Ergebnis eindeutig?
\item \textbf{Kontext:} Waren Zielgruppe, Daten und Grenzen vorhanden?
\item \textbf{Werkzeug:} Wurde die richtige Fähigkeit mit passenden Parametern verwendet?
\item \textbf{Evidenz:} Sind Quellen, Berechnungen und Ableitungen tragfähig?
\item \textbf{Ausgabe:} Wurde das Ergebnis inhaltlich und visuell geprüft?
\end{enumerate}

\section{Die kleinste wirksame Korrektur}

Ändere bei einem fehlgeschlagenen Versuch möglichst nur eine Variable. Ergänze beispielsweise zuerst die Zielgruppe, dann separat ein Format. Werden gleichzeitig zehn Regeln geändert, bleibt unklar, welche davon die Verbesserung verursacht hat. Dieses Vorgehen macht Prompt- und Workflow-Entwicklung zu einem nachvollziehbaren Experiment.

\section{Wann man nicht weiter automatisieren sollte}

Ein Workflow ist ein schlechter Automatisierungskandidat, wenn der Zweck bei jedem Durchlauf wechselt, zentrale Daten nicht zuverlässig verfügbar sind oder eine kleine Fehlentscheidung große Schäden verursachen kann. In solchen Fällen ist KI als Analyse- oder Entwurfswerkzeug oft wertvoller als als autonomer Ausführer.

\begin{praxis}
Nimm ein enttäuschendes Ergebnis aus einem früheren Chat. Ordne den Fehler einer der fünf Diagnosefragen zu und ändere nur diese Ebene. Dokumentiere, ob die zweite Ausgabe messbar besser geworden ist.
\end{praxis}



\appendix
\chapter{Cheatsheet}
\begin{lstlisting}
/analyse +tief       Ursachen und Zusammenhänge
/entscheidung +kritisch +tabelle
/recherche +quellen  aktuelle, belegte Recherche
/plan +schritt       ausführbarer Ablauf
/dokument +datei     fertiges Artefakt
\end{lstlisting}
\textbf{Auftragsformel:} Ergebnis + Kontext + Grenzen + Qualitätskriterien + gewünschtes Format.

\chapter{Quellen und Stand}
Stand der produktbezogenen Hinweise: 8. August 2026. Oberflächen und Verfügbarkeit können sich ändern.
\begin{itemize}
\item OpenAI Help Center: \href{https://help.openai.com/en/articles/10500283-deep-research}{Deep Research in ChatGPT}.
\item OpenAI Help Center: \href{https://help.openai.com/en/articles/11487775-connectors-in}{Apps in ChatGPT}.
\item OpenAI Help Center: \href{https://help.openai.com/en/articles/8096356-custom-instructions-for-chatgpt}{Custom Instructions}.
\item OpenAI Help Center: \href{https://help.openai.com/en/articles/8554407}{GPTs in ChatGPT}.
\end{itemize}

\end{document}
